% ---- pg-042 （书页 30） ----

(a) 为所有变量作边缘点图.

(b) 构造阵列 $\bar{\bm{x}}, \bm{S}$, 和 $\bm{R}$ 并对 $\bm{R}$ 中元素作解释.

\mvbold{1.7} 下面给出对 $p = 2$ 个变量的 $n = 3$ 个观测值：

变量 1：\quad $x_{11} = 2$ \quad $x_{21} = 3$ \quad $x_{31} = 4$

变量 2：\quad $x_{12} = 1$ \quad $x_{22} = 2$ \quad $x_{32} = 4$

(a) 为每对观测值在二维“变量空间”作图，即构造数据的二维散布图.

(b) 在三维“项目空间”中将数据作为两点标出.

\mvbold{1.8} 用\hyperref[eq:1-12]{式(1-12)} 的欧几里德距离公式，取 $p = 2$ 及用\hyperref[eq:1-20]{式(1-20)} 的统计距离公式，取 $a_{11} = 1/3, a_{22} = 4/27, a_{12} = 1/9$. 计算点 $P = (-1, 1)$ 到点 $Q = (1, 0)$ 的距离. 画出到点 $Q$ 的统计距离平方为常数 1 的点的轨迹.

\mvbold{1.9} 考虑下面两个变量 $x_1$ 和 $x_2$ 的 8 对测量值：

\begin{mvtable}
\begin{tabular}{lcccccccc}
\toprule
$x_1$ & $-6$ & $-3$ & $-2$ & 1 & 2 & 5 & 6 & 8 \\
$x_2$ & $-2$ & $-3$ & 1 & $-1$ & 2 & 1 & 5 & 3 \\
\bottomrule
\end{tabular}
\end{mvtable}

(a) 画出数据的散布图，并计算 $s_{11}, s_{22}$ 及 $s_{12}$.

(b) 假定将原始坐标轴旋转角度 $\theta = 26^\circ$ [已知 $\cos(26^\circ) = 0.899$ 和 $\sin(26^\circ) = 0.438$]. 利用\hyperref[eq:1-18]{式(1-18)} 计算与变量 $x_1, \tilde{x}_2$ 对应的测量值.

(c) 根据 (b) 中求出的 $\tilde{x}_1, \tilde{x}_2$ 的测量值，计算样本方差 $\tilde{s}_{11}$ 及 $\tilde{s}_{22}$.

(d) 设一对新的测量值 $(x_1, x_2) = (4, -2)$. 用\hyperref[eq:1-18]{式(1-18)} 将它们变换为 $\tilde{x}_1$ 和 $\tilde{x}_2$ 的测量值，并用\hyperref[eq:1-17]{式(1-17)} 计算原点 $O = (0, 0)$ 到新点 $P = (\tilde{x}_1, \tilde{x}_2)$ 的距离 $d(O, P)$.

注意：需要用到 (c) 中求出的 $\tilde{s}_{11}$ 和 $\tilde{s}_{22}$.

(e) 用\hyperref[eq:1-19]{式(1-19)} 及脚注 2 中给出的关于 $a_{11}, a_{12}$ 及 $a_{22}$ 的表达式，计算点 $P = (4, -2)$ 到原点 $O = (0, 0)$ 的距离.

注意：需用到 (a) 中求出的 $s_{11}, s_{22}$ 及 $s_{12}$.

将这一计算出的距离和 (d) 中用 $\tilde{x}_1$ 和 $\tilde{x}_2$ 的值算出的距离进行比较（在舍入误差范围内，数值应是相同的）.

\mvbold{1.10} 下面的距离函数对于计算从原点出发的距离是否有效？试解释.

(a) $x_1^2 + 4x_2^2 + x_1x_2 = $(距离)$^2$

(b) $x_1^2 - 2x_2^2 = $(距离)$^2$

\mvbold{1.11} 设 $a_{11} = 4, a_{22} = 1$ 和 $a_{12} = -1$，证明\hyperref[eq:1-20]{式(1-20)} 定义的距离满足\hyperref[eq:1-25]{式(1-25)} 中的前三个条件（三角不等式较难证明）.

\mvbold{1.12} 定义从点 $P = (x_1, x_2)$ 到原点 $O = (0, 0)$ 的距离为
\[
d(O, P) = \max\{|x_1|, |x_2|\}
\]

(a) 计算从点 $P = (-3, 4)$ 到原点的距离.

(b) 画出到原点的距离平方为 1 的点的轨迹.

(c) 将上述距离公式推广到 $p$ 维空间中的点.

\mvbold{1.13} 一个大城市的主要道路被设计成网格形，如下图所示. 街道 1 至街道 5 的方向为从北至南 (NS)，街道 A 至街道 E 方向为从东至西 (EW). 假设在交点 $(A, 2), (E, 3)$ 和 $(C, 5)$ 上有零售商店. 又假定在一条街上两个交点之间的距离不论是南北方向还是东西方向都为 1 个单位. 网格上任何两个交点（点）间的距离称为“城市街区”距离. （例如交点 $(D, 1)$ 和 $(C, 2)$ 间的距离，我们可以称之为 $d((D, 1), (C, 2))$，可以这样算出 $d((D, 1), (C, 2)) = d((D, 1)$，
